\section{The Correspondence Project}

Beginning tonight, on on successive Mondays, we plan to make available selected correspondence between the various men of tradition. Our initial efforts will include letters from \textbf{Rene Guenon} to \textbf{Guido de Giorgio} (the letters in the opposite direction are missing) as well as the exchanges between \textbf{Mircea Eliade} and \textbf{Julius Evola}. We will learn, for example, that Guenon was in touch with the three advocates of Romanity, albeit with quite different approaches: De Giorgio, Evola, and Reghini. The letters to De Giorio are instructive since they often offer critiques of the other two.

I know from the statistics that our texts by De Giorgio are among the least read. That is unfortunate for, as we shall see, Guenon, unlike the other two, treats him as an equal, a peer. Tradition is much more than discussing what some king may have said to a priest three millennia ago; it is nothing if it is not accompanied by a change of consciousness, the recovery and re-creation of Traditional modes of thinking and experiencing the world in one's very being. De Giorgio accomplished that, he internalized Tradition. He did not always explicitly ``spell it out'', but it is obvious to the discerning reader. Besides absorbing the two Traditions of Rome, he brings in elements of the Vedanta and Sufism into his meditations.

The exchanges between Guenon and De Giorgio take place between 1925 and and 1930, just prior to Guenon's relocation to Cairo, Egypt. We can therefore account for Guenon's whereabouts during this period. We will see his true opinion of Evola. We will learn of a Sufi center for initiation in Paris, but not in Rome. This should silence our critics who claim we just fabricate facts. Gornahoor is a center for cooperation, and those who only understand \emph{strife} place themselves outside. They are simply fighting the wrong battle, but we pray they will eventually engage in the greater and lesser Holy Wars.


\flrightit{Posted on 2012-06-18 by Admin}

\begin{center}* * *\end{center}

\begin{footnotesize}
\begin{sffamily}
\texttt{Justin on 2012-06-18 at 08:27 said:}

A tantalising prospect!

\hfill

\texttt{golgonooza on 2012-06-18 at 12:29 said:}

On the subject of things to read, and knowing that you intend this summer to read The Divine Comedy as an initiatic text -- I was thinking of reading it (only in translation unfortunately) for the first time but wanted to be have some kind of esoteric guide (aside from Virgil!) -- is there any commentary on it you recommend? (Sorry this is a bit off-topic)

\hfill

\texttt{Cologero on 2012-06-18 at 18:34 said:}

The first goal is to understand the literal meaning of the text which is hardly possible today because of all the allusions to ancient mythology, poetry, and historical events. I use the Charles Singleton prose translation and commentary. The penguin translation by Musa also has excellent notes.

As for an ``esoteric'' guide, try Guenon's book on Dante. Dante starts by indicating he is taking the journey of ``our'' life, so he is inviting us to participate, not to be mere observers or curiosity seekers. He makes the journey as a ``living man'', so we shouldn't read it as a ``hell and brimstone'' sermon about our future fate. Rather, we need to look at the spiritual tendencies he describes in our own consciousness. ``Fear'' holds us back, since we don't want to face such things; after all, we tend to think more highly of ourselves than is warranted. Least of all, should we read it as a theological treatise, describing God's justice, etc., as is commonly done. Nor is it a travelogue or a vision of the fate of specific persons, who really represent ``types'' or symbols, of certain states of consciousness.

It must also be read in the light of the goals of the ``Fedeli d'Amore'', those men faithful to love (of which ``sex magick'' is a corruption). This tradition is mostly lost in the West, so some familiarity with the Sufi version (e.g., Rumi, Hafiz, Ibn Arabi) may be helpful.

\end{sffamily}
\end{footnotesize}

